\documentclass[12pt]{article}
\usepackage[brazil]{babel}
\usepackage{amsmath,amsthm}
\newcommand{\R}{\mathbb{R}}
\newcommand{\destaque}[1]{\textbf{#1}}
\DeclareMathOperator{\sen}{sen}
\newtheorem{teorema}{Teorema}
\newtheorem{definicao}{Definição}
\title{Um estudo sobre funções}
\author{Ana Souza \and Bruno Lima}
\date{Novembro de 2026}

\begin{document}
\maketitle

\begin{abstract}
Este trabalho revisa \emph{funções reais} e mostra o uso de \LaTeX{} no dia a dia.
\end{abstract}

\section{Introdução}\label{sec:intro}
Seja $f\colon \R \to \R$ uma função\footnote{Contínua, para simplificar.}. Usamos o
\destaque{teorema} da Seção~\ref{sec:teo} e a Equação~\eqref{eq:euler}. % comentário que some
O custo foi de R\$~10 e 50\% do tempo --- segundo \cite{knuth}.

\subsection{Notação}
\begin{itemize}
  \item Conjunto dos reais: $\R$;
  \item Seno: $\sen x$;
  \item[$\boxtimes$] tarefa feita
  \begin{enumerate}
    \item primeiro
    \item segundo
  \end{enumerate}
\end{itemize}

\begin{equation}\label{eq:euler}
  e^{i\pi} + 1 = 0
\end{equation}

\begin{align}
  a &= b + c \label{eq:a}\\
  d &= e \nonumber
\end{align}

\section{Resultados}\label{sec:teo}

\begin{definicao}
Uma função é \textit{par} se $f(-x) = f(x)$.
\end{definicao}

\begin{teorema}[Pitágoras]
Em um triângulo retângulo, $a^2 = b^2 + c^2$.
\end{teorema}

\begin{proof}
Óbvio.
\end{proof}

\begin{table}[h]
\centering
\caption{Derivadas}\label{tab:der}
\begin{tabular}{|l|c|}
\hline
Função & Derivada \\ \hline
$x^n$ & $n x^{n-1}$ \\
$|x|$ & $\operatorname{sgn}(x)$ \\
\multicolumn{2}{c}{fim} \\
\hline
\end{tabular}
\end{table}

Veja a Tabela~\ref{tab:der}.

\begin{lstlisting}[language=Python]
def f(x):  # 100% {ok}
    return x**2
\end{lstlisting}

Código inline: \verb|a_b%c| e \texttt{x\_y}. Link: \href{https://example.com}{site} e \url{https://a.org/x%20y}.
Química: \ce{H2O}.

\begin{quote}
Uma citação.
\end{quote}

\begin{thebibliography}{9}
\bibitem{knuth} D. Knuth. \emph{The \TeX book}. 1984.
\end{thebibliography}
\end{document}
